\medskip

On étudie l'évolution de la population d'une espèce animale au sein d'une réserve naturelle.
Les effectifs de cette population ont été recensés à différentes années. Les données collectées sont présentées dans le tableau suivant :

\begin{center}
	\begin{tabularx}{0.9\linewidth}{|l|*{4}{>{\centering \arraybackslash}X|}} \hline
		Année              & 2000 & 2005 & 2010 & 2015 \\ \hline
		Nombre d'individus & 50   & 64   & 80   & 100  \\	\hline
	\end{tabularx}
\end{center}


Pour anticiper l'évolution de cette population, la direction de la réserve a choisi de modéliser le nombre d'individus en fonction du temps.
Pour cela, elle utilise une fonction, définie sur l'intervalle $[0 ~;~ +\infty [$, dont la variable $x$ représente le temps écoulé, en année, à partir de l'année 2000.
Dans son modèle, l'image de 0 par cette fonction vaut 50, ce qui correspond au nombre d'individus en l'an 2000.

	\subsection*{Partie A. Modèle 1}

	Dans cette partie, la direction de la réserve fait l'hypothèse que la fonction cherchée satisfait l'équation différentielle suivante :
$y'=0,05 y-0,5 \quad\left(E_{1}\right)$

	\begin{enumerate}
		\item 
\begin{list}{\textbullet}{On résout l'équation différentielle $\left(E_{1}\right)$ avec la condition initiale $y(0)=50$.}
\item L'équation différentielle $y=ay$ a pour solutions les fonctions $y$ vérifiant\\ $y(x)=C\e^{ax}$ où $C$ est un réel quelconque.

Donc l'équation différentielle $y=0,05 y$ a pour solutions les fonctions $y$ vérifiant $y(x)=C\e^{0,05 x}$ où $C$ est un réel quelconque.

\item On cherche une solution particulière constante $k$ à l'équation $\left (E_1\right )$.

On a donc $0=0,05k-0,5$ donc $k=10$.
\item Les solutions de l'équation $\left (E_1\right )$ sont donc les fonctions $y$ définies par\\ $y(x)=C\e^{-0,05x}+10$ où $C\in\R$.
\item On cherche la solution vérifiant la condition initiale $y(0)=50$.

On a donc $C\e^{0}+10=50$ donc $C=40$.
\end{list}

La solution de l'équation différentielle $\left (E_1\right )$ vérifiant $y(0)=50$ est la fonction $y$ définie par $y(x)=40\e^{0,05x}+10$.

		\item On compare les résultats du tableau avec ceux que l'on obtiendrait avec ce modèle.
		
\begin{center}
	\begin{tabularx}{0.9\linewidth}{|c|*{4}{>{\centering \arraybackslash}X|}} \hline
		Année              & 2000 & 2005 & 2010 & 2015 \\ \hline
		Nombre d'individus & 50   & 64   & 80   & 100  \\	\hline
		$x$ & 0 & 5 & 10 & 15 \\ \hline
		$y(x)$ & $50$ & $61,36$ & $75,95$ & $94,68$ \\ \hline
	\end{tabularx}
\end{center}		
		
	\end{enumerate}

	\subsection*{Partie B. Modèle 2}

	Dans cette partie, la direction de la réserve fait l'hypothèse que la fonction cherchée satisfait l'équation différentielle suivante :
$y'=0,05 y \left (1-0,00125 y\right ) \quad (E_2)$.

On note $f$ la fonction définie sur $[0 ~;~ +\infty[$ par :
$f(x)=\dfrac{800}{1+15 \e^{-0,05 x}}$,

	et $\mathcal{C}$ sa courbe représentative dans un repère orthonormé.

	\medskip

	À l'aide d'un logiciel de calcul formel, on a obtenu les résultats suivants.



	\begin{center}
\renewcommand{\arraystretch}{1.5}
\begin{tabularx}{\linewidth}{|l|l|>{\centering \arraybackslash}X|} 
\hline
  & Instruction & Résultat \\	
\hline
1 & $f(x):=\dfrac{800}{1+15 \e^{-0,05 x}}$	& $f(x)=\dfrac{800}{1+15 \e^{-0,05 x}}$  \rule[-13pt]{0pt}{30pt}\\	
\hline
2 & $f'(x):=$ Dérivée $(f(x))$ & $f'(x)=\dfrac{600 \e^{-0,05 x}}{\left(1+15 \e^{-0,05 x}\right)^{2}}$ \rule[-15pt]{0pt}{35pt} \\ 
\hline
3 & $f''(x):=$ Dérivée ( $f'(x)$ ) & $f''(x)=\dfrac{30 \e^{-0,05 x}}{\left(1+15 \e^{-0,05 x}\right)^{3}}\left(15 \e^{-0,05 x}-1\right)$ \rule[-15pt]{0pt}{35pt} \\ 
\hline
4 & Résoudre ( $15 \e^{-0,05 x}-1 \geqslant 0$ )  & $x \leqslant 20 \ln (15)$  \\		
\hline
\end{tabularx}
\renewcommand{\arraystretch}{1}
	\end{center}

	\begin{enumerate}
		\item 
On a d'une part:		$f(0)=\dfrac{800}{1+15 \e^{0}} = \dfrac{800}{16}=50$
		
\begin{list}{\textbullet}{D'autre part:}
\item $0,05f(x)= \dfrac{0,05\times 800}{1+15\e^{-0,05x}} = \dfrac{40}{1+15\e^{-0,05x}}$
\item $\np{0,00125} f(x) = \dfrac{\np{0,00125}\times 800}{1+15\e^{-0,05x}}= \dfrac{1}{1+15\e^{-0,05x}}$
\item $1- \np{0,00125} f(x) = 1- \dfrac{1}{1+15\e^{-0,05x}}= \dfrac{1+15\e^{-0,05x}-1}{1+15 \e^{-0,05x}} = \dfrac{15 \e^{-0,05x}}{1+15 \e^{-0,05x}}$
\item 
$\begin{aligned}[t]
0,05f(x) \left (1- \np{0,00125} f(x)  \right ) 
& = \dfrac{40}{1+15\e^{-0,05x}} \times \dfrac{15 \e^{-0,05x}}{1+15 \e^{-0,05x}} 
= \dfrac{40\times 15 \e^{-0,05x}}{\left (1+15 \e^{-0,05x}\right )^2} \\
& = \dfrac{600 \e^{-0,05x}}{\left (1+15 \e^{-0,05x}\right )^2}
= f'(x) \text{ d'après le tableau}
\end{aligned}$
\end{list}

On admet que cette fonction $f$ est l'unique solution de $(E_{2})$ prenant la valeur initiale de 50 en 0.

		\item Avec ce nouveau modèle $f$,  l'effectif de cette population en 2050 peut être estimé à $f(50)$ soit environ 358.
		


		\item 
$\ds\lim_{x\to +\infty} -0,05x = -\infty$ et
$\ds\lim_{X\to -\infty} \e^{X}=0$ donc 
$\ds\lim_{x\to +\infty} \e^{-0,05x} = 0$

On en déduit  que 
$\ds\lim_{x\to +\infty} 1+15\e^{-0,05x} = 1$ et que
$\ds\lim_{x\to +\infty} \dfrac{800}{1+15\e^{-0,05x}} = 800$.

Donc $\ds\lim_{x\to +\infty} f(x)=800$, ce qui prouve que la courbe $\mathcal{C}$ admet la droite d'équation $y=800$ comme asymptote horizontale en $+\infty$.
		
		\item 
Sur $[0 ~;~+\infty[$,  $f'(x)=\dfrac{600 \e^{-0,05 x}}{\left(1+15 \e^{-0,05 x}\right)^{2}}$.

Or pour tout réel $X$, on sait que $\e^{X}>0$, donc $\e^{-0,05x}>0$ pour tout $x$.

$\left (1+15\e^{-0,05x}\right )^2>0$ donc $f'(x)>0$.

On en conclut que la fonction $f$ est croissante sur $[0 ~;~+\infty[$.

		\item 
On résout l'inéquation $15 \e^{-0,05 x}-1 \geqslant 0$.

$\begin{aligned}
15 \e^{-0,05 x}-1 \geqslant 0
& \iff 15 \e^{-0,05 x} \geqslant 1
 \iff  \e^{-0,05 x} \geqslant \dfrac{1}{15}
 \iff  -0,05 x \geqslant \ln\left (\dfrac{1}{15}\right )\\
& \iff  x \leqslant \dfrac{ \ln\left (\frac{1}{15}\right )}{-0,05}
 \iff  x \leqslant \dfrac{ -\ln\left (15 \right )}{-0,05}
 \iff  x \leqslant 20\,\ln\left (15 \right )
\end{aligned}$

		\item On admet que la vitesse de croissance de la population de cette espèce, exprimée en nombre d'individus par an, est modélisée par la fonction $f'$.
		\begin{enumerate}
			\item
Pour étudier la convexité de la fonction $f$ sur l'intervalle $[0 ~;~ +\infty[$, on détermine le signe de la dérivée seconde $f''(x)$.

D'après le tableau: $f''(x)=\dfrac{30 \e^{-0,05 x}}{\left(1+15 \e^{-0,05 x}\right)^{3}}\left(15 \e^{-0,05 x}-1\right)$.

D'après ce qui a été vu précédemment, $f''(x)$ est du signe de $\left(15 \e^{-0,05 x}-1\right)$.

\[ \begin{tablvar}[8em]{2}
 \hline
 x & 0 & & 20\,\ln(15) & & +\infty\\
 \hline
15 \e^{-0,05 x}-1 & & + & \barre[0] & - & \\
 \hline
f''(x) & & + & \barre[0] & - & \\
 \hline
  & & f \text{ convexe} & \barre & f \text{ concave} & \\
 \hline
 \end{tablvar}\]
 
En $x=20\,\ln(15)$, $f''(x)$ s'annule et change de signe, donc la courbe $\mathcal{C}$ admet un point d'inflexion.

Si $x=20\,\ln(15)$, $15\e^{-0,05x}-1=0$ donc $15\e^{-0,05\times 20\ln(15)}=1$

$f\left (20\ln(15)\right )
= \dfrac{800}{1+ 15\e^{-0,05\times 20\ln(15)}}
= \dfrac{800}{1+1} = 400$

Donc le point de coordonnées $(20\,\ln(15)\,,\, 400)$ est le seul point d'inflexion de la courbe $\mathcal{C}$.
 
			\item La direction de la réserve affirme :
			\og Au vu de ce modèle, la vitesse de croissance de la population de cette espèce va augmenter pendant un peu plus de cinquante ans, puis va diminuer \fg{}. 
			


On établit le tableau de variations de la fonction $f'$ qui modélise  la vitesse de croissance de la population de cette espèce, exprimée en nombre d'individus par an.

\[ \begin{tablvar}[8em]{2}
 \hline
 x & 0 & & 20\,\ln(15) & & +\infty\\
 \hline
f''(x) & & + & \barre[0] & - & \\
 \hline
\variations{\mil{f'} & \bas{} && \haut{} &&  \bas{}} 
\hline
\end{tablvar}\]

Or $20\,\ln(15) \approx 54,2$ donc la direction a raison.
		\end{enumerate}
	\end{enumerate}

\bigskip
